\begin{figure}[H]
    \centering
    \begin{tikzpicture}[
        node distance=1.15cm and 1.25cm,
        every node/.style={font=\small}
    ]
        \node[tesisService, minimum width=4.0cm] (app) {Aplicación Spring Boot};
        \node[tesisProcess, below=of app, minimum width=4.7cm] (lib) {Librería Transactional Outbox};

        \node[tesisBox, fill=tesisSoftBlue, below left=1.25cm and 2.05cm of lib, minimum width=2.6cm] (persist) {Módulo\\Persistencia};
        \node[tesisBox, fill=tesisSoftAmber, below=1.25cm of lib, minimum width=2.6cm] (publisher) {Módulo\\Publicador};
        \node[tesisBox, fill=tesisSoftGreen, below right=1.25cm and 2.05cm of lib, minimum width=2.6cm] (idempotent) {Módulo\\Idempotencia};

        \node[tesisDatabase, below=of persist, minimum width=2.6cm] (db) {PostgreSQL};
        \node[tesisBroker, below=of publisher, minimum width=2.6cm] (kafka) {Apache Kafka};
        \node[tesisDatabase, below=of idempotent, minimum width=2.6cm] (processed) {Eventos\\procesados};

        \draw[tesisArrow] (app) -- (lib);
        \draw[tesisArrow] (lib) -- (persist);
        \draw[tesisArrow] (lib) -- (publisher);
        \draw[tesisArrow] (lib) -- (idempotent);
        \draw[tesisArrow] (persist) -- (db);
        \draw[tesisArrow] (publisher) -- (kafka);
        \draw[tesisDashedArrow] (idempotent) -- (processed);
    \end{tikzpicture}

    \caption{Arquitectura modular de la librería propuesta.}
    \label{fig:arquitectura-libreria}
\end{figure}
